\documentclass{subfiles}
\begin{document}
    \begin{definition*}
        A finite set of symbols \(A\) is said to be an alphabet,
            if its elements are used to represet strings.
    \end{definition*}

    \begin{definition*}
        Let \(A\) be an alphabet. We define word (or string) an \(n\)-tuple of 
            \(A^{n}\). That is, 
            \[
                A^{n} = \set*{\bv{w} = (x_{1}, \ldots, x_{n})}[x_{i} \in A, \forall i].
            \]
        Additionally, given a word \(\bv{w}\) its weight \(\norm{\bv{w}}\) is the number of
            non-zero components in \(\bv{w}\).
    \end{definition*}

    From the above definition two remarks can be made about the weight of a word:
        \begin{enumerate}
            \item For any alphabet \(A, \bv{w} \in A^{n}\) has weight 0 if and only if 
                \(\bv{w} = \bv{0}\).

            \item For any two strings \(\bv{w}_{1}, \bv{w}_{2} \in A^{n}\), it holds
                \[
                    \norm{\bv{w}_{1} - \bv{w}_{2}} = 0 \iff \bv{w}_{1} = \bv{w}_{2}.
                \]
        \end{enumerate}

    \begin{definition*}[Hamming's distance]
        Let \(\bv{w}_{1}, \bv{w}_{2} \in A^{n}\), for some alphabet \(A\).
            We define the (Hamming) distance of \(\bv{w}_{1}, \bv{w}_{2}\)
            to be the number of components for which these differ.
    \end{definition*}
    \begin{remark*}
        It can be easely proved that the Hamming's distance is a metric.
    \end{remark*}

    \begin{definition*}[Code]
        Let \(A\) be an alphabet, and consider the vector space \(A^{n}\) these creates.
            We define code a subset \(\C*\) of \(A^{n}\).
    \end{definition*}

    \begin{definition*}[Minimal distance]
        Let \(\C*\) be a code. We say that \(C*\) has Minimal distance \(d\)
        if 
        \[
            d = \min{\abs{\bv{w}}}[\bv{w} \in \C*].
        \]
    \end{definition*}

    That is, the minimal distance of a code is the minimal weight of all the words in \(\c*\).

    In our discussion of codes we assume \(A\) to be always some finite field \(\F\),
    or equivalently a Galoi field of \(q\) elements \(\gf{q}\).
    \clearpage
\end{document}
